\emph{This chapter was written by Claude, Anthropic's AI, from an
outline and direction by Carsten Wulff, who reviewed and edited the
result. The commit history of the book's repository records precisely
who wrote what.}

Analog design that cannot be reproduced from a shell command is not
done. That is the thesis of this chapter, and the reason the tools in it
exist. Every schematic in this course simulates from a script, every
testbench sweeps its corners from a one-line command, and one of the
ADCs in the converter chapter compiles its own layout. None of that
needs a licence server.

The tools share a prefix - cic, for Custom IC Creator - and a
philosophy: plain text in, plain text out, and git in between. Each one
is small. Together they are a design flow.

\section{cicsim - simulations you can
rerun}\label{cicsim---simulations-you-can-rerun}

\texttt{cicsim} is the workhorse: a script package that controls
ngspice. You met it in the tutorial when \texttt{cicsim\ simcell} built
the simulation directory for the current mirror. What it buys you over
running ngspice by hand:

\texttt{cicsim\ run} takes a testbench and a corner specification -
typical, slow, fast, high and low temperature, supply corners - and runs
the cross product, in parallel if you ask (\texttt{-\/-threads}), with a
progress bar and a timeout to kill the simulation that wedged. Add
\texttt{-\/-count\ N} and each corner runs N times with mismatch - Monte
Carlo from the same command line.

The measurements come back as a results table, one row per corner, which
is the artifact that matters: the OTA chapter's advice to verify gain,
noise, slew and start-up \emph{over PVT} is one \texttt{cicsim\ run} per
testbench, and the table is the evidence.

The yaml files that define a simulation live next to the schematic and
go into git. Six months later, \texttt{cicsim\ run} again and you get
the same table - or a diff that tells you exactly what the PDK update
broke.

\begin{itemize}
\tightlist
\item
  \texttt{cicsim\ run} : corners, in parallel, with a results table
\item
  \texttt{-\/-count\ N} : Monte Carlo from the same command
\item
  yaml + git : the simulation is reproducible, or it is not done
\end{itemize}

\section{ciccreator - the layout
compiler}\label{ciccreator---the-layout-compiler}

\texttt{ciccreator} is the tool behind the compiled SAR ADC in the
converter chapter. The idea dates to 2013: describe the circuit as a
SPICE netlist, describe the layout intent as a JSON object definition,
add a technology rule file, and let a compiler place the polygons. The
prototype was 16 thousand lines of Perl; the current tool is a C++
rewrite of the same input language, fast enough to recompile an ADC
while you watch.

The point is portability. The netlist and the object file do not mention
a technology; the rule file does. Porting the ADC to a new process - and
it has been ported to 22, 28, 55, 65, 130 and 180 nm - means writing a
new rule file, not redrawing a layout. That is why the comparison table
in the converter chapter has a ``compiled'' line with a competitive
figure of merit, and why the same ADC exists as an open source SKY130
implementation you can read.

\begin{itemize}
\tightlist
\item
  SPICE netlist + JSON object definition + technology rules = layout
\item
  Port to a new process: new rule file, not new polygons
\item
  The compiled SAR ADC of the converter chapter is the proof
\end{itemize}

\section{cicpy and cicconf - the
glue}\label{cicpy-and-cicconf---the-glue}

\texttt{cicpy} translates. It takes ciccreator output and transpiles it
to whatever the rest of your flow speaks: SKILL for Cadence, SPICE,
Verilog, Xschem schematics, Magic layout, or SVG when you just want to
look at a cell. It also carries two small workhorses - \texttt{sch2mag}
and \texttt{spi2mag} - that netlist a schematic and place-and-route it
to Magic, which is how the standard-cell-like blocks in the aicex IPs
appear.

\texttt{cicconf} manages the projects themselves. An aicex IP is a git
repository, and a chip is many of them; \texttt{cicconf} keeps the
dependency list in one \texttt{config.yaml} and clones, updates and
templates the collection. When the tutorial told you to run
\texttt{cicconf\ clone}, this is what stitched your project together.

\begin{itemize}
\tightlist
\item
  cicpy: transpile the compiled IC to SKILL, SPICE, Verilog, Xschem,
  Magic or SVG
\item
  cicconf: one config.yaml for a project of many git repositories
\end{itemize}

\section{cicwave - looking at the
waves}\label{cicwave---looking-at-the-waves}

\texttt{cicwave} is the waveform viewer: it reads ngspice raw files and
draws them with a PyQtGraph/Qt6 backend fast enough for long transients.
It is the third viewer of its line - the first was written during a
summer internship in 2001 - and it exists for the same reason as the
rest of the family: the open source flow deserves a viewer that starts
instantly, works on every OS, and is scriptable.

\section{cictikz - the figures of this
book}\label{cictikz---the-figures-of-this-book}

The newest member drew the book you are reading. \texttt{cictikz}
packages the TikZ symbol library behind every schematic figure in these
chapters, gives an AI assistant a render-and-look loop so figures can be
drawn and reviewed programmatically, and converts between the book's
TikZ dialect and Xschem schematics - so a figure can start life as a
real, simulated schematic and end as a book drawing, or the other way
around.

\begin{itemize}
\tightlist
\item
  One symbol library for every schematic in the book
\item
  TikZ to Xschem and back: figures that are also schematics
\end{itemize}

\section{Summary}\label{summary}

The one-page version of this chapter:

\begin{itemize}
\tightlist
\item
  If it does not rerun from a shell command, it is not done
\item
  cicsim: corners and Monte Carlo as one command, results as a table,
  all of it in git
\item
  ciccreator: netlist + object definition + rules compile to layout;
  porting is a rule file
\item
  cicpy transpiles to the flow you have; cicconf holds multi-repo
  projects together
\item
  cicwave views the waves; cictikz draws the book
\item
  None of it needs a licence server
\end{itemize}

\section{Would you like to know
more?}\label{would-you-like-to-know-more}

Every tool lives on GitHub with its own documentation:
\href{https://github.com/wulffern/cicsim}{cicsim},
\href{https://github.com/wulffern/ciccreator}{ciccreator}
(\href{https://ciccreator.readthedocs.io/en/latest/index.html}{docs}),
\href{https://github.com/wulffern/cicpy}{cicpy},
\href{https://github.com/wulffern/cicconf}{cicconf},
\href{https://github.com/wulffern/cicwave}{cicwave}
(\href{https://wulffern.github.io/cicwave/}{docs}) and
\href{https://github.com/wulffern/cictikz}{cictikz}
(\href{https://analogicus.com/cictikz/}{docs}). The IPs built with them
are collected in \href{https://github.com/wulffern/aicex}{aicex}.

\begin{itemize}
\tightlist
\item
  \href{https://github.com/wulffern/aicex}{github.com/wulffern/aicex} -
  the IPs built with these tools
\end{itemize}
